% ============================================================
\section{Ưu điểm, hạn chế và hướng mở rộng}
\label{sec:discussion}
% ============================================================

\Cref{tab:proscons} tóm tắt ưu điểm và hạn chế theo từng khía cạnh.

\begin{table}[H]
  \centering
  \small
  \begin{tabularx}{\linewidth}{|P{2.2cm}|L|L|}
    \hline
    \textbf{Khía cạnh} & \textbf{Ưu điểm} & \textbf{Hạn chế} \\ \hline
    Văn phạm & Mọi câu trả lời lần ngược được về cây cú pháp & Độ phủ đóng: câu đúng miền
      nhưng diễn đạt lạ bị từ chối, như C007, C008 \\ \hline
    Parser & Earley chạy trên mọi CFG, kể cả luật rỗng và đệ quy trái & Chỉ phát hiện mơ hồ
      tới hai cây \\ \hline
    Trả lời & Không bịa dữ kiện; ba loại thất bại được báo riêng & Chủ đề, tài liệu tra theo
      từ khóa; KB thiếu quan hệ tài liệu--chương \\ \hline
    Ngữ cảnh & Giải đúng 6/6 hội thoại mẫu & Chỉ năm dạng tham chiếu; chưa hỏi lại khi có
      nhiều ứng viên \\ \hline
    Vận hành & Chạy offline, không cần model hay GPU, kết quả tất định & -- \\ \hline
    Kiểm thử & Từng tầng được kiểm thử độc lập & Chưa có tập đánh giá độc lập, nên 96,88\%
      chưa đại diện cho câu hỏi thật \\ \hline
  \end{tabularx}
  \caption{Ưu điểm và hạn chế}
  \label{tab:proscons}
\end{table}

Hướng mở rộng tiếp theo là thu thập câu hỏi thật của sinh viên làm tập đánh giá độc lập,
thêm số thứ tự và từ đồng nghĩa vào văn phạm, thử PCFG hoặc so khớp mềm cho câu chỉ lệch văn
phạm ít, bổ sung dữ kiện có nguồn vào KB như giảng viên và ngày thi, và hỏi lại người dùng
khi tham chiếu có nhiều ứng viên.
